\label{kap:komplexni}
\mathquote{The shortest path between two truths in the real domain passes through the complex domain.}{Jacques Hadamard}

V kapitole 2 jsme slíbili, že se ke komplexním číslům vrátíme. V kapitole~\ref{kap:kvadraticke} jsme viděli, že rovnice $x^2 + 2x + 5 = 0$ nemá reálné řešení, a v příběhu o Cardanovi, že se odmocninám ze záporných čísel nedalo vyhnout ani při řešení rovnic, které reálné řešení mají. Komplexní čísla dlouho měla pověst něčeho podezřelého. Descartes je pohrdavě nazval \textit{imaginárními}, tedy pomyslnými, a to jméno jim už zůstalo. Dnes bez nich nefunguje elektrotechnika, kvantová fyzika ani zpracování signálů. Jsou totiž jediným rozumným způsobem, jak počítat s otáčením.

\section{Algebraický tvar}\label{sec:algebraicky-tvar}

\begin{definitionbox}
  \textbf{Imaginární jednotka} $i$ je číslo, pro které platí $i^2 = -1$.

  \textbf{Komplexní číslo} je číslo tvaru
  \begin{equation*}
    z = a + bi,\qquad a, b\in\mathbb{R}.
  \end{equation*}
  Číslo $a = \mathrm{Re}\,z$ je \textbf{reálná část}, $b = \mathrm{Im}\,z$ \textbf{imaginární část}. Množinu všech komplexních čísel značíme $\mathbb{C}$. Pro $b = 0$ dostáváme reálná čísla, takže $\mathbb{R}\subset\mathbb{C}$. Čísla s $a = 0$ se nazývají \textbf{ryze imaginární}.

  Dvě komplexní čísla se rovnají, právě když se rovnají jejich reálné i imaginární části.
\end{definitionbox}

V kapitole 2 jsme psali $i = \sqrt{-1}$. To je dobrá intuice, ale přesnější je definovat $i$ rovností $i^2 = -1$. Odmocnina ze záporného čísla totiž v $\mathbb{C}$ není jednoznačná (číslo $-i$ splňuje ${(-i)}^2 = -1$ úplně stejně) a pravidla jako $\sqrt{a}\sqrt{b} = \sqrt{ab}$ přestávají platit: $\sqrt{-1}\cdot\sqrt{-1} = i\cdot i = -1$, ale $\sqrt{(-1)(-1)} = \sqrt{1} = 1$.

\begin{definitionbox}
  \textbf{Operace}: Pro $z_1 = a + bi$ a $z_2 = c + di$:
  \begin{align*}
    z_1 + z_2 &= (a + c) + (b + d)i,\\
    z_1\cdot z_2 &= (ac - bd) + (ad + bc)i.
  \end{align*}
  Násobíme tedy jako mnohočleny a pak nahradíme $i^2 = -1$.

  \textbf{Komplexně sdružené číslo} k $z = a + bi$ je $\bar{z} = a - bi$. Platí $z\cdot\bar{z} = a^2 + b^2$ (reálné číslo!).

  \textbf{Dělení}: zlomek rozšíříme číslem komplexně sdruženým se jmenovatelem:
  \begin{equation*}
    \frac{z_1}{z_2} = \frac{z_1\cdot\bar{z}_2}{z_2\cdot\bar{z}_2} = \frac{z_1\cdot\bar{z}_2}{c^2 + d^2}.
  \end{equation*}
\end{definitionbox}

Dělení je přesně stejný trik jako usměrňování zlomků v kapitole~\ref{kap:vyrazy}. Tam jsme se zbavovali odmocniny ve jmenovateli pomocí $(a + \sqrt{b})(a - \sqrt{b}) = a^2 - b$, tady se zbavujeme $i$ pomocí $(c + di)(c - di) = c^2 + d^2$.

\begin{example}
  Pro $z_1 = 2 + 3i$ a $z_2 = 1 - i$:
  \begin{align*}
    z_1 + z_2 &= 3 + 2i,\\
    z_1\cdot z_2 &= 2 - 2i + 3i - 3i^2 = 2 + i + 3 = 5 + i,\\
    \frac{z_1}{z_2} &= \frac{(2 + 3i)(1 + i)}{(1 - i)(1 + i)} = \frac{2 + 2i + 3i + 3i^2}{1 + 1} = \frac{-1 + 5i}{2} = -\frac{1}{2} + \frac{5}{2}i.
  \end{align*}
\end{example}

\begin{note}
  Mocniny $i$ se opakují s periodou $4$:
  \begin{equation*}
    i^1 = i,\quad i^2 = -1,\quad i^3 = -i,\quad i^4 = 1,\quad i^5 = i,\ \ldots
  \end{equation*}
  Stačí tedy vydělit exponent čtyřmi a podívat se na zbytek. Například $i^{2026} = i^{4\cdot 506 + 2} = i^2 = -1$.
  Proč zrovna perioda $4$? Za chvíli v Gaussově rovině uvidíme, že násobení číslem $i$ znamená otočení o $90^\circ$. Každá další mocnina je tedy o čtvrt otáčky dál a po čtyřech čtvrtotáčkách jsme zpátky. Vpravo je totéž otáčení použité na číslo $z = 2 + i$: čísla $z, iz, i^2z, i^3z$ tvoří čtverec.
  \begin{center}
    \begin{tikzpicture}[scale=0.95]
      \begin{scope}
        \draw[->] (-1.7,0) -- (1.8,0) node[right] {\tiny Re};
        \draw[->] (0,-1.6) -- (0,1.7) node[above] {\tiny Im};
        \draw[dashed, gray] (0,0) circle (1.2);
        \foreach \a/\l/\p in {0/{$i^0 = 1$}/above left, 90/{$i$}/below right, 180/{$i^2 = -1$}/below right, 270/{$i^3 = -i$}/above left} {
          \filldraw[colorofwisdom] (\a:1.2) circle (2pt) node[\p, font=\scriptsize] {\l};
        }
        \foreach \a in {0,90,180,270} \draw[->, colorofwisdom, thick] (\a+12:1.45) arc (\a+12:\a+78:1.45);
        \node[font=\scriptsize, colorofwisdom] at (45:1.8) {$\cdot i$};
      \end{scope}
      \begin{scope}[xshift=5.5cm, scale=0.6]
        \draw[->] (-2.8,0) -- (2.9,0) node[right] {\tiny Re};
        \draw[->] (0,-2.7) -- (0,2.8) node[above] {\tiny Im};
        \draw[gray, dashed] (2,1) -- (-1,2) -- (-2,-1) -- (1,-2) -- cycle;
        \foreach \x/\y/\l/\p in {2/1/{$z = 2 + i$}/right, -1/2/{$iz = -1 + 2i$}/above left, -2/-1/{$i^2z$}/left, 1/-2/{$i^3z$}/below right} {
          \draw[colorofwisdom] (0,0) -- (\x,\y);
          \filldraw[colorofwisdom] (\x,\y) circle (3pt) node[\p, font=\scriptsize] {\l};
        }
      \end{scope}
    \end{tikzpicture}
  \end{center}
\end{note}

\begin{note}\label{note:komplexni-usporadani}
  Komplexní čísla tvoří \textbf{těleso} z kapitoly 2: platí všechny axiomy sčítání a násobení a každé nenulové číslo má inverzi $\frac{1}{z} = \frac{\bar{z}}{|z|^2}$. Na rozdíl od reálných čísel ale \textbf{nejsou uspořádaná}. Nedává smysl říct, že $i > 0$ nebo $i < 0$. Kdyby totiž $i > 0$, pak $i\cdot i > 0$, tedy $-1 > 0$, a kdyby $i < 0$, pak $(-i)(-i) > 0$, tedy zase $-1 > 0$. Nerovnice v komplexních číslech proto řešit nemůžeme.
\end{note}

\begin{problem}
  Vypočítejte:
  \begin{enumerate}[label=(\alph*)]
    \item $(3 - 2i)(4 + i)$,
    \item ${(1 + i)}^2$ a ${(1 + i)}^4$,
    \item $\dfrac{5}{2 - i}$,
    \item $\dfrac{1 + i}{1 - i}$,
    \item $i + i^2 + i^3 + \cdots + i^{100}$.
  \end{enumerate}
\end{problem}

\begin{solution}
  \begin{enumerate}[label=(\alph*)]
    \item $12 + 3i - 8i - 2i^2 = 14 - 5i$.
    \item ${(1 + i)}^2 = 1 + 2i - 1 = 2i$, ${(1 + i)}^4 = {(2i)}^2 = -4$.
    \item $\frac{5(2 + i)}{4 + 1} = 2 + i$.
    \item $\frac{{(1 + i)}^2}{2} = \frac{2i}{2} = i$.
    \item Každá čtveřice $i + (-1) + (-i) + 1 = 0$, a čtveřic je $25$. Součet je $0$.
  \end{enumerate}
\end{solution}

\section{Gaussova rovina}\label{sec:gaussova-rovina}

Komplexní číslo $a + bi$ je určeno dvojicí reálných čísel $(a, b)$. Můžeme ho tedy zakreslit jako bod v rovině. Vodorovná osa je \textbf{reálná}, svislá \textbf{imaginární}. Takové rovině se říká \textbf{Gaussova} (nebo komplexní) rovina.

\begin{definitionbox}
  \textbf{Absolutní hodnota} komplexního čísla $z = a + bi$ je jeho vzdálenost od počátku:
  \begin{equation*}
    |z| = \sqrt{a^2 + b^2} = \sqrt{z\bar{z}}.
  \end{equation*}
  \textbf{Argument} $\arg z$ je orientovaný úhel, který svírá spojnice počátku s bodem $z$ s kladnou reálnou poloosou.

  \vspace{0.2cm}

  \textbf{Geometrický význam operací}:
  \begin{itemize}
    \item sčítání komplexních čísel je sčítání vektorů z kapitoly~\ref{kap:analyticka},
    \item komplexně sdružené číslo je obraz v osové souměrnosti podle reálné osy,
    \item $|z_1 - z_2|$ je vzdálenost bodů $z_1$ a $z_2$.
  \end{itemize}
\end{definitionbox}

\begin{center}
  \begin{tikzpicture}[scale=1.1]
    \draw[->] (-1,0) -- (4.2,0) node[right] {\small Re};
    \draw[->] (0,-2.5) -- (0,3) node[above] {\small Im};
    \coordinate (Z) at (3,2);
    \draw[thick, colorofwisdom, ->] (0,0) -- (Z);
    \filldraw[colorofwisdom] (Z) circle (1.5pt) node[above right] {\small $z = a + bi$};
    \filldraw (3,-2) circle (1.5pt) node[below right] {\small $\bar{z} = a - bi$};
    \draw[dashed, gray] (3,-2) -- (Z);
    \draw[dashed] (Z) -- (3,0) node[below right] {\small $a$};
    \draw[dashed] (Z) -- (0,2) node[left] {\small $b$};
    \node[colorofwisdom, above left] at (1.5,1) {\small $|z|$};
    \draw[->] (0.6,0) arc (0:33.7:0.6);
    \node at (0.85,0.25) {\small $\varphi$};
  \end{tikzpicture}
\end{center}

Komplexní číslo si tedy můžeme představit dvojím způsobem: jako \textbf{bod} $[a, b]$, nebo jako \textbf{šipku} (vektor) z počátku do tohoto bodu. Sčítání je pak skládání šipek, přesně jako u sil ve fyzice. Šipku $z_2$ přiložíme na konec šipky $z_1$, a kde skončíme, tam je $z_1 + z_2$. Obě šipky a jejich součet tvoří rovnoběžník. Na obrázku jsou čísla z příkladu výše, $z_1 = 2 + 3i$ a $z_2 = 1 - i$.

\begin{center}
  \begin{tikzpicture}[scale=0.85]
    \draw[->] (-0.8,0) -- (4.3,0) node[right] {\small Re};
    \draw[->] (0,-1.6) -- (0,3.7) node[above] {\small Im};
    \foreach \x in {1,2,3} \draw (\x,0.06) -- (\x,-0.06);
    \foreach \y in {-1,1,2,3} \draw (0.06,\y) -- (-0.06,\y);
    \fill[colorofwisdom!8] (0,0) -- (2,3) -- (3,2) -- (1,-1) -- cycle;
    \draw[->, thick, colorofwisdom] (0,0) -- (2,3) node[above left, font=\small] {$z_1 = 2 + 3i$};
    \draw[->, thick] (0,0) -- (1,-1) node[below right, font=\small] {$z_2 = 1 - i$};
    \draw[->, dashed] (2,3) -- (3,2);
    \draw[->, dashed, colorofwisdom] (1,-1) -- (3,2);
    \draw[->, very thick, colorofwisdom!60!black] (0,0) -- (3,2) node[right, font=\small] {$z_1 + z_2 = 3 + 2i$};
  \end{tikzpicture}
  \hspace{1cm}
  \begin{tikzpicture}[scale=0.85]
    \draw[->] (-0.8,0) -- (4.3,0) node[right] {\small Re};
    \draw[->] (0,-1.6) -- (0,3.7) node[above] {\small Im};
    \draw[->, thick, colorofwisdom] (0,0) -- (2,3) node[above left, font=\small] {$z_1$};
    \draw[->, thick] (0,0) -- (1,-1) node[below right, font=\small] {$z_2$};
    \draw[very thick, colorofwisdom!60!black] (1,-1) -- (2,3) node[midway, right, font=\small] {$|z_1 - z_2| = \sqrt{17}$};
    \draw[dashed, gray] (0,0) circle (1.414);
    \node[gray, font=\scriptsize] at (-0.3,-1.65) {$|z| = \sqrt{2}$};
  \end{tikzpicture}
\end{center}
Vpravo je vzdálenost dvou čísel: $z_1 - z_2 = 1 + 4i$, a tedy $|z_1 - z_2| = \sqrt{1 + 16} = \sqrt{17}$. Šedá kružnice je množina všech čísel se stejnou absolutní hodnotou jako $z_2$, tedy $\{z : |z| = \sqrt{2}\}$. Rovnice $|z - z_0| = r$ obecně popisuje kružnici se středem $z_0$ a poloměrem $r$.

\section{Goniometrický tvar}\label{sec:goniometricky-tvar}

\begin{definitionbox}
  Každé nenulové komplexní číslo lze zapsat v \textbf{goniometrickém tvaru}
  \begin{equation*}
    z = |z|(\cos\varphi + i\sin\varphi),
  \end{equation*}
  kde $\varphi = \arg z$. Pro $z = a + bi$ platí
  \begin{equation*}
    \cos\varphi = \frac{a}{|z|},\qquad\sin\varphi = \frac{b}{|z|}.
  \end{equation*}

  \textbf{Násobení a dělení v goniometrickém tvaru}: Pro $z_1 = r_1(\cos\varphi_1 + i\sin\varphi_1)$ a $z_2 = r_2(\cos\varphi_2 + i\sin\varphi_2)$:
  \begin{align*}
    z_1\cdot z_2 &= r_1r_2\left(\cos(\varphi_1 + \varphi_2) + i\sin(\varphi_1 + \varphi_2)\right),\\
    \frac{z_1}{z_2} &= \frac{r_1}{r_2}\left(\cos(\varphi_1 - \varphi_2) + i\sin(\varphi_1 - \varphi_2)\right).
  \end{align*}
  Při násobení se tedy \textbf{absolutní hodnoty násobí} a \textbf{argumenty sčítají}.
\end{definitionbox}

\begin{proof}[(násobení)]
  Roznásobíme a použijeme součtové vzorce z kapitoly~\ref{kap:goniometrie}:
  \begin{align*}
    z_1z_2 &= r_1r_2\left(\cos\varphi_1\cos\varphi_2 + i\cos\varphi_1\sin\varphi_2 + i\sin\varphi_1\cos\varphi_2 + i^2\sin\varphi_1\sin\varphi_2\right)\\
    &= r_1r_2\left((\cos\varphi_1\cos\varphi_2 - \sin\varphi_1\sin\varphi_2) + i(\sin\varphi_1\cos\varphi_2 + \cos\varphi_1\sin\varphi_2)\right)\\
    &= r_1r_2\left(\cos(\varphi_1 + \varphi_2) + i\sin(\varphi_1 + \varphi_2)\right).
  \end{align*}
\end{proof}

Tohle je ten nejdůležitější poznatek celé kapitoly. Násobení komplexním číslem s absolutní hodnotou $1$ je \textbf{otočení} roviny. Násobení číslem $i$ (které má $|i| = 1$ a $\arg i = \frac{\pi}{2}$) otočí každý bod o $90^\circ$ proti směru hodinových ručiček. Dvojí otočení o $90^\circ$ je otočení o $180^\circ$, tedy násobení $-1$. A to je geometrický důvod, proč $i^2 = -1$! Nic mystického, jen dvě otočení za sebou.

\begin{center}
  \begin{tikzpicture}[scale=1.15]
    \draw[->] (-1.4,0) -- (2.6,0) node[right] {\small Re};
    \draw[->] (0,-0.3) -- (0,3.1) node[above] {\small Im};
    \coordinate (Z1) at (1,1);
    \coordinate (Z2) at (1,1.732);
    \coordinate (P) at (-0.732,2.732);
    \draw[->, thick, gray] (0,0) -- (Z1) node[right, font=\small, black] {$z_1 = 1 + i$};
    \draw[->, thick, black] (0,0) -- (Z2) node[right, font=\small] {$z_2 = 1 + \sqrt{3}\,i$};
    \draw[->, very thick, colorofwisdom] (0,0) -- (P) node[above left, font=\small] {$z_1z_2$};
    \draw[gray] (0.55,0) arc (0:45:0.55);
    \node[font=\scriptsize, gray] at (22:0.75) {$45^\circ$};
    \draw (60:0.95) arc (60:105:0.95);
    \node[font=\scriptsize] at (83:1.15) {$45^\circ$};
    \draw (1.25,0) arc (0:60:1.25);
    \node[font=\scriptsize] at (30:1.48) {$60^\circ$};
    \draw[->, colorofwisdom] (1.8,0) arc (0:105:1.8);
    \node[font=\scriptsize, colorofwisdom] at (12:2.1) {$105^\circ$};
    \node[anchor=west, align=left, font=\small] at (3,1.6) {$|z_1| = \sqrt{2}$, $\arg z_1 = 45^\circ$\\$|z_2| = 2$, $\arg z_2 = 60^\circ$\\[2pt]$|z_1z_2| = 2\sqrt{2}$, $\arg z_1z_2 = 105^\circ$\\[2pt]$z_1z_2 = (1 - \sqrt{3}) + (1 + \sqrt{3})i$};
  \end{tikzpicture}
\end{center}
Na obrázku násobíme $z_2 = 1 + \sqrt{3}\,i$ číslem $z_1 = 1 + i$. Šipka $z_2$ se otočí o $\arg z_1 = 45^\circ$ a natáhne se $|z_1| = \sqrt{2}$-krát. Zkuste si ověřit roznásobením, že $(1 + i)(1 + \sqrt{3}\,i)$ vyjde opravdu tak, jak je napsáno vpravo. Násobení komplexních čísel je tedy vždy kombinace dvou jednoduchých pohybů: \textbf{otočení a natažení}.

\begin{example}
  Převeďme do goniometrického tvaru:
  \begin{itemize}
    \item $z = 1 + i$: $|z| = \sqrt{2}$, $\cos\varphi = \sin\varphi = \frac{1}{\sqrt{2}}$, tedy $\varphi = \frac{\pi}{4}$ a $z = \sqrt{2}\left(\cos\frac{\pi}{4} + i\sin\frac{\pi}{4}\right)$.
    \item $z = -1 + i\sqrt{3}$: $|z| = 2$, $\cos\varphi = -\frac{1}{2}$, $\sin\varphi = \frac{\sqrt{3}}{2}$, tedy $\varphi = \frac{2\pi}{3}$ (druhý kvadrant!).
    \item $z = -5$: $|z| = 5$, $\varphi = \pi$, $z = 5(\cos\pi + i\sin\pi)$.
    \item $z = -3i$: $|z| = 3$, $\varphi = \frac{3\pi}{2}$.
  \end{itemize}
  Při určování argumentu se vždy podívejte, ve kterém kvadrantu číslo leží. Samotný tangens nestačí, protože $\tg\varphi = \frac{b}{a}$ je stejný pro $z$ i $-z$.
\end{example}

\begin{definitionbox}\label{def:moivre}
  \textbf{Moivreova věta}: Pro každé $n\in\mathbb{Z}$ platí
  \begin{equation*}
    {\left(r(\cos\varphi + i\sin\varphi)\right)}^n = r^n(\cos n\varphi + i\sin n\varphi).
  \end{equation*}
\end{definitionbox}

\begin{proof}
  Pro přirozená $n$ matematickou indukcí. Pro $n = 1$ je to zřejmé. Pokud tvrzení platí pro $n$, pak podle pravidla o násobení
  \begin{equation*}
    z^{n+1} = z^n\cdot z = r^n(\cos n\varphi + i\sin n\varphi)\cdot r(\cos\varphi + i\sin\varphi) = r^{n+1}\left(\cos(n + 1)\varphi + i\sin(n + 1)\varphi\right).
  \end{equation*}
  Pro záporná $n$ to plyne z pravidla pro dělení.
\end{proof}

Moivreova věta říká graficky, že mocniny $z, z^2, z^3, \ldots$ leží na \textbf{spirále}. Každé další vynásobení číslem $z$ otočí o stejný úhel $\varphi$ a natáhne $|z|$-krát. Pro $|z| > 1$ se spirála rozvíjí ven, pro $|z| < 1$ se stáčí do počátku a pro $|z| = 1$ body obíhají po jednotkové kružnici. Na obrázku jsou mocniny čísla $1 + i$, každá otočená o $45^\circ$ a $\sqrt{2}$-krát delší než předchozí.

\begin{center}
  \begin{tikzpicture}[scale=0.6]
    \draw[->] (-5,0) -- (2.6,0) node[right] {\small Re};
    \draw[->] (0,-4.6) -- (0,2.9) node[above] {\small Im};
    \draw[thick, colorofwisdom!50, domain=0:225, samples=80] plot ({pow(2,\x/90)*cos(\x)}, {pow(2,\x/90)*sin(\x)});
    \foreach \x/\y/\l/\p in {1/0/{$z^0 = 1$}/below right, 1/1/{$z = 1 + i$}/right, 0/2/{$z^2 = 2i$}/above right, -2/2/{$z^3 = -2 + 2i$}/above left, -4/0/{$z^4 = -4$}/above left, -4/-4/{$z^5 = -4 - 4i$}/below right} {
      \draw[gray] (0,0) -- (\x,\y);
      \filldraw[colorofwisdom] (\x,\y) circle (3.5pt) node[\p, font=\scriptsize] {\l};
    }
  \end{tikzpicture}
\end{center}

\begin{example}
  \begin{itemize}
    \item ${(1 + i)}^8 = {\left(\sqrt{2}\right)}^8\left(\cos 2\pi + i\sin 2\pi\right) = 16$. Zkuste to roznásobit osmkrát.
    \item ${(1 + i\sqrt{3})}^6 = 2^6(\cos 2\pi + i\sin 2\pi) = 64$, protože $1 + i\sqrt{3} = 2\left(\cos\frac{\pi}{3} + i\sin\frac{\pi}{3}\right)$.
    \item Moivreova věta dá zadarmo vzorce pro násobné úhly. Pro $n = 2$ a $r = 1$: ${(\cos\varphi + i\sin\varphi)}^2 = \cos^2\varphi - \sin^2\varphi + 2i\sin\varphi\cos\varphi$, a to se musí rovnat $\cos 2\varphi + i\sin 2\varphi$. Porovnáním reálných a imaginárních částí dostaneme oba vzorce pro dvojnásobný úhel z kapitoly~\ref{kap:goniometrie}.
  \end{itemize}
\end{example}

\begin{problem}
  \begin{enumerate}[label=(\alph*)]
    \item Převeďte do goniometrického tvaru $\sqrt{3} - i$ a $-2 - 2i$.
    \item Vypočítejte ${(\sqrt{3} - i)}^{12}$.
    \item Vypočítejte $\dfrac{{(1 + i)}^{10}}{{(1 - i)}^{8}}$.
  \end{enumerate}
\end{problem}

\begin{solution}
  \begin{enumerate}[label=(\alph*)]
    \item $\sqrt{3} - i = 2\left(\cos\frac{11\pi}{6} + i\sin\frac{11\pi}{6}\right)$, $-2 - 2i = 2\sqrt{2}\left(\cos\frac{5\pi}{4} + i\sin\frac{5\pi}{4}\right)$.
    \item $2^{12}\left(\cos 22\pi + i\sin 22\pi\right) = 4096$.
    \item $1 + i = \sqrt{2}\left(\cos\frac{\pi}{4} + i\sin\frac{\pi}{4}\right)$ a $1 - i = \sqrt{2}\left(\cos\left(-\frac{\pi}{4}\right) + i\sin\left(-\frac{\pi}{4}\right)\right)$. Podíl má absolutní hodnotu ${\left(\sqrt{2}\right)}^{10 - 8} = 2$ a argument $10\cdot\frac{\pi}{4} - 8\cdot\left(-\frac{\pi}{4}\right) = \frac{9\pi}{2}$, tedy výsledek je $2\left(\cos\frac{9\pi}{2} + i\sin\frac{9\pi}{2}\right) = 2i$.
  \end{enumerate}
\end{solution}

\section{Rovnice v komplexních číslech}

\subsection{Kvadratické rovnice}\label{sec:komplexni-kvadraticke}

\begin{definitionbox}
  Kvadratická rovnice $ax^2 + bx + c = 0$ s reálnými koeficienty a $D < 0$ má v $\mathbb{C}$ dva \textbf{komplexně sdružené} kořeny
  \begin{equation*}
    x_{1,2} = \frac{-b\pm i\sqrt{-D}}{2a}.
  \end{equation*}
  Vietovy vzorce i rozklad $a(x - x_1)(x - x_2)$ platí dál.
\end{definitionbox}

\begin{example}
  Rovnice $x^2 + 2x + 5 = 0$ z kapitoly~\ref{kap:kvadraticke}: $D = -16$, $\sqrt{-D} = 4$, tedy $x_{1,2} = \frac{-2\pm 4i}{2} = -1\pm 2i$. Kontrola Vietem: součet $-2$, součin ${(-1)}^2 + 2^2 = 5$. Sedí.

  \begin{center}
    \begin{tikzpicture}[scale=0.8]
      \begin{scope}
        \draw[->] (-3.4,0) -- (1.6,0) node[right] {\small $x$};
        \draw[->] (0,-0.4) -- (0,6) node[above] {\small $y$};
        \draw[thick, colorofwisdom, domain=-3.2:1.2, samples=60] plot (\x, {(\x*\x + 2*\x + 5)/1.6});
        \filldraw (-1,2.5) circle (2pt) node[below, font=\scriptsize] {$[-1, 4]$};
        \node[font=\small] at (-1,-1) {$y = x^2 + 2x + 5$};
      \end{scope}
      \begin{scope}[xshift=6.5cm, yshift=2.6cm]
        \draw[->] (-2.4,0) -- (1.6,0) node[right] {\small Re};
        \draw[->] (0,-2.6) -- (0,2.7) node[above] {\small Im};
        \draw[dashed, gray] (-1,2) -- (-1,-2);
        \filldraw[colorofwisdom] (-1,2) circle (2.5pt) node[left, font=\scriptsize] {$-1 + 2i$};
        \filldraw[colorofwisdom] (-1,-2) circle (2.5pt) node[left, font=\scriptsize] {$-1 - 2i$};
        \draw (-1,0.06) -- (-1,-0.06) node[below right, font=\scriptsize] {$-1$};
        \node[font=\small] at (-0.4,-3.6) {kořeny v Gaussově rovině};
      \end{scope}
    \end{tikzpicture}
  \end{center}
  Parabola osu $x$ neprotíná, takže reálné kořeny nejsou. Komplexní kořeny leží symetricky podle reálné osy (jsou komplexně sdružené) a jejich reálná část $-1$ je přesně $x$-ová souřadnice vrcholu paraboly. To platí obecně: $-\frac{b}{2a}$ je střed mezi kořeny, ať jsou reálné, nebo komplexní.
\end{example}

\begin{definitionbox}
  \textbf{Základní věta algebry}: Každý mnohočlen stupně $n\geq 1$ s komplexními koeficienty má v $\mathbb{C}$ alespoň jeden kořen. Důsledek: lze ho rozložit na součin $n$ lineárních činitelů, takže má (počítáno s násobností) právě $n$ kořenů.
\end{definitionbox}

Tuto větu dokázal (ne úplně kompletně) Carl Friedrich Gauss ve své disertaci v roce 1799, když mu bylo dvaadvacet. Je to důvod, proč u komplexních čísel nemusíme vymýšlet žádná další čísla. U reálných jsme museli přidat $i$, aby měla řešení rovnice $x^2 + 1 = 0$. U komplexních čísel už mají řešení všechny polynomiální rovnice. Komplexní čísla jsou v tomto smyslu `hotová'.

\subsection{Binomické rovnice}\label{sec:binomicke-rovnice}

\begin{definitionbox}
  \textbf{Binomická rovnice} $z^n = a$, kde $a = r(\cos\alpha + i\sin\alpha)\neq 0$, má v $\mathbb{C}$ právě $n$ řešení
  \begin{equation*}
    z_k = \sqrt[n]{r}\left(\cos\frac{\alpha + 2k\pi}{n} + i\sin\frac{\alpha + 2k\pi}{n}\right),\qquad k = 0, 1, \ldots, n - 1.
  \end{equation*}
  V Gaussově rovině tvoří tato řešení vrcholy \textbf{pravidelného $n$-úhelníku} se středem v počátku.
\end{definitionbox}

Proč $+2k\pi$? Podle Moivreovy věty má $z^n$ argument $n\varphi$. Ten se musí rovnat $\alpha$, ale jen až na násobek $2\pi$, protože úhly lišící se o $2\pi$ jsou stejné. Tedy $n\varphi = \alpha + 2k\pi$. Pro $k = 0, \ldots, n - 1$ dostaneme různé body, pro $k = n$ se vrátíme na začátek.

\begin{example}
  Rovnice $z^3 = 1$ má řešení $\cos\frac{2k\pi}{3} + i\sin\frac{2k\pi}{3}$ pro $k = 0, 1, 2$:
  \begin{equation*}
    z_0 = 1,\qquad z_1 = -\frac{1}{2} + \frac{\sqrt{3}}{2}i,\qquad z_2 = -\frac{1}{2} - \frac{\sqrt{3}}{2}i.
  \end{equation*}
  Tvoří rovnostranný trojúhelník vepsaný do jednotkové kružnice. Číslo $1$ má tedy tři třetí odmocniny, jen jedna z nich je reálná.

  Rovnice $z^4 = -16$: $r = 16$, $\alpha = \pi$, tedy $z_k = 2\left(\cos\frac{\pi + 2k\pi}{4} + i\sin\frac{\pi + 2k\pi}{4}\right)$. Úhly $\frac{\pi}{4}, \frac{3\pi}{4}, \frac{5\pi}{4}, \frac{7\pi}{4}$ dávají čtyři vrcholy čtverce $\pm\sqrt{2}\pm\sqrt{2}i$.
  \begin{center}
    \begin{tikzpicture}[scale=1.1]
      \begin{scope}
        \draw[->] (-1.5,0) -- (1.6,0) node[right] {\tiny Re};
        \draw[->] (0,-1.4) -- (0,1.5) node[above] {\tiny Im};
        \draw[dashed, gray] (0,0) circle (1);
        \draw[thick, colorofwisdom] (0:1) -- (120:1) -- (240:1) -- cycle;
        \foreach \a in {0,120,240} \filldraw[colorofwisdom] (\a:1) circle (1.5pt);
        \node at (0,-1.8) {\small $z^3 = 1$};
      \end{scope}
      \begin{scope}[xshift=5cm]
        \draw[->] (-1.8,0) -- (1.9,0) node[right] {\tiny Re};
        \draw[->] (0,-1.7) -- (0,1.8) node[above] {\tiny Im};
        \draw[dashed, gray] (0,0) circle (1.3);
        \draw[thick, colorofwisdom] (45:1.3) -- (135:1.3) -- (225:1.3) -- (315:1.3) -- cycle;
        \foreach \a in {45,135,225,315} \filldraw[colorofwisdom] (\a:1.3) circle (1.5pt);
        \node at (0,-1.8) {\small $z^4 = -16$};
      \end{scope}
    \end{tikzpicture}
  \end{center}
\end{example}

\begin{problem}
  Vyřešte v $\mathbb{C}$:
  \begin{enumerate}[label=(\alph*)]
    \item $x^2 - 6x + 13 = 0$,
    \item $x^2 + 9 = 0$,
    \item $z^2 = i$,
    \item $z^6 = 64$,
    \item $x^3 + 8 = 0$.
  \end{enumerate}
\end{problem}

\begin{solution}
  \begin{enumerate}[label=(\alph*)]
    \item $D = 36 - 52 = -16$, $x = 3\pm 2i$.
    \item $x = \pm 3i$.
    \item $i = \cos\frac{\pi}{2} + i\sin\frac{\pi}{2}$, řešení $\cos\frac{\pi}{4} + i\sin\frac{\pi}{4} = \frac{\sqrt{2}}{2}(1 + i)$ a $\cos\frac{5\pi}{4} + i\sin\frac{5\pi}{4} = -\frac{\sqrt{2}}{2}(1 + i)$.
    \item $z_k = 2\left(\cos\frac{k\pi}{3} + i\sin\frac{k\pi}{3}\right)$ pro $k = 0, \ldots, 5$: $\pm 2$, $\pm 1\pm\sqrt{3}i$. Pravidelný šestiúhelník.
    \item $x^3 + 8 = (x + 2)(x^2 - 2x + 4)$ podle vzorce z kapitoly~\ref{kap:vyrazy}. Kořeny $-2$ a $1\pm\sqrt{3}i$.
  \end{enumerate}
\end{solution}

\section{Příběh --- Muž, který předpověděl svou smrt}

Abraham de Moivre, autor Moivreovy věty, byl francouzský protestant. Když v roce 1685 Ludvík XIV. zrušil edikt nantský, který protestantům zaručoval náboženskou svobodu, byl de Moivre uvězněn a teprve po propuštění mohl uprchnout do Anglie. Tam strávil zbytek života. Navzdory tomu, že se přátelil s Newtonem a Halleym a byl členem Královské společnosti, nikdy nezískal univerzitní místo. Živil se doučováním matematiky, poradenstvím pro hazardní hráče a pojišťovny a tím, že v londýnské kavárně Slaughter's hrál o peníze šachy.

V pozdním věku si prý všiml, že každý den spí o čtvrt hodiny déle než předchozí den. Jako správný matematik z toho sestavil aritmetickou posloupnost z kapitoly~\ref{kap:posloupnosti} a spočítal, kdy dosáhne dvaceti čtyř hodin. Toho dne, 27. listopadu 1754, opravdu zemřel. Lékař jako příčinu smrti uvedl somnolenci, tedy chorobnou spavost. Je to pravděpodobně jediný případ v dějinách, kdy někdo předpověděl datum své smrti pomocí aritmetické posloupnosti a vyšlo mu to. Matematika je někdy až příliš přesná.

A ještě jedna věc, kterou nemůžeme vynechat. Pokud se v goniometrickém tvaru podíváme na $\cos\varphi + i\sin\varphi$, Euler v 18. století ukázal, že se to rovná $e^{i\varphi}$, kde $e$ je Eulerovo číslo z kapitoly~\ref{kap:exponencialni}. Moivreova věta je pak jen pravidlo ${(e^{i\varphi})}^n = e^{in\varphi}$. A pro $\varphi = \pi$ dostaneme
\begin{equation*}
  e^{i\pi} + 1 = 0,
\end{equation*}
rovnici, která v jednom řádku spojuje pět nejdůležitějších čísel matematiky: $0$, $1$, $\pi$, $e$ a $i$. Mnoho matematiků ji považuje za nejkrásnější rovnici, jaká kdy byla napsána. Fyzik Richard Feynman ji ve svém deníku ve čtrnácti letech nazval `nejpozoruhodnějším vzorcem v matematice'.

\section{Shrnutí a otázky k ústní zkoušce}

\begin{problem}
  Odpovězte vlastními slovy:
  \begin{enumerate}
    \item Proč potřebujeme komplexní čísla? Jak je definována imaginární jednotka?
    \item Co je algebraický tvar komplexního čísla? Jak se komplexní čísla sčítají, násobí a dělí?
    \item Co je komplexně sdružené číslo a absolutní hodnota? Jaký mají geometrický význam?
    \item Proč komplexní čísla nelze uspořádat?
    \item Co je Gaussova rovina? Jak v ní vypadá sčítání?
    \item Co je goniometrický tvar? Jak se v něm násobí a dělí? Dokažte to.
    \item Vyslovte a dokažte Moivreovu větu.
    \item Jak se řeší kvadratické rovnice se záporným diskriminantem? Co říká základní věta algebry?
    \item Jak se řeší binomické rovnice? Proč jejich řešení tvoří pravidelný mnohoúhelník?
  \end{enumerate}
\end{problem}
